\message{ !name(GSoSD_foundational principle developments.tex)}\documentclass[a4paper]{arrowhead}

\usepackage[yyyymmdd]{datetime}
\usepackage{etoolbox}
\usepackage[utf8]{inputenc}
\usepackage{multirow}

\renewcommand{\dateseparator}{-}

%% Special references
\newcommand{\fref}[1]{{\textcolor{ArrowheadBlue}{\hyperref[sec:functions:#1]{#1}}}}
\newcommand{\mref}[1]{{\textcolor{ArrowheadPurple}{\hyperref[sec:model:#1]{#1}}}}
\newcommand{\pdef}[1]{{\textcolor{ArrowheadGrey}{#1\label{sec:model:primitives:#1}\label{sec:model:primitives:#1s}\label{sec:model:primitives:#1es}}}}
\newcommand{\pref}[1]{{\textcolor{ArrowheadGrey}{\hyperref[sec:model:primitives:#1]{#1}}}}

\newrobustcmd\fsubsection[3]{
  \addtocounter{subsection}{1}
  \addcontentsline{toc}{subsection}{\protect\numberline{\thesubsection}function \textcolor{ArrowheadBlue}{#1}}
  \renewcommand*{\do}[1]{\rref{##1},\ }
  \subsection*{
    \thesubsection\quad
    operation
    \textcolor{ArrowheadBlue}{#1}
    (\notblank{#2}{\mref{#2}}{})
    \notblank{#3}{: \mref{#3}}{}
  }
  \label{sec:functions:#1}
}
\newrobustcmd\msubsection[2]{
  \addtocounter{subsection}{1}
  \addcontentsline{toc}{subsection}{\protect\numberline{\thesubsection}#1 \textcolor{ArrowheadPurple}{#2}}
  \subsection*{\thesubsection\quad#1 \textcolor{ArrowheadPurple}{#2}}
  \label{sec:model:#2} \label{sec:model:#2s} \label{sec:model:#2es}
}
%%

\begin{document}

\message{ !name(GSoSD_foundational principle developments.tex) !offset(-3) }


%% Arrowhead Document Properties
\ArrowheadTitle{XXX} % XXX = SystemName e.g. Service Registry HTTP/TLS/JSON}
\ArrowheadType{System Description}
\ArrowheadTypeShort{SysD}
\ArrowheadVersion{X.Y.Z} % Arrowhead version X.Y.Z, e..g. 4.4.1
\ArrowheadDate{\today}
\ArrowheadAuthor{Doument author} % Corresponding author e.g. Jerker Delsing
\ArrowheadStatus{RELEASE} % e..g. RELEASE, RELEASE CONDIDATE, PROTOTYPE
\ArrowheadContact{xx@yy.xx} % Email of corresponding author
\ArrowheadFooter{\href{www.arrowhead.eu}{www.arrowhead.eu}}
\ArrowheadSetup
%%

%% Front Page
\begin{center}
  \vspace*{1cm}
  \huge{\arrowtitle}

  \vspace*{0.2cm}
  \LARGE{\arrowtype}
  \vspace*{1cm}

  %\Large{Service ID: \textit{"\arrowid"}}
  \vspace*{\fill}

  % Front Page Image
  %\includegraphics{figures/TODO}

  \vspace*{1cm}
  \vspace*{\fill}

  % Front Page Abstract
  \begin{abstract}
    This is the template for System Description (SysD document)
    according to the Eclipse Arrowehad documentation structure. 
  \end{abstract}

  \vspace*{1cm}

%   \scriptsize
%   \begin{tabularx}{\textwidth}{l X}
%     \raisebox{-0.5\height}{\includegraphics[width=2cm]{figures/artemis_logo}} & {ARTEMIS Innovation Pilot Project: Arrowhead\newline
%     THEME [SP1-JTI-ARTEMIS-2012-AIPP4 SP1-JTI-ARTEMIS-2012-AIPP6]\newline
%     [Production and Energy System Automation Intelligent-Built environment and urban infrastructure for sustainable and friendly cities]}
%   \end{tabularx}
%   \vspace*{-0.2cm}
 \end{center}

\newpage
%%

%% Table of Contents
\tableofcontents
\newpage
%%

\section{Overview}
\label{sec:overview}
This document describes the fundational principels for Eclispe
Arrowhead. 

The rest of this document is organized as follows.
In Section \ref{sec:prior_art}, we reference major prior art capabilitites
of the system.
In Section \ref{sec:use}, we the intended usage of the system.
In Section \ref{sec:properties}, we describe fundmental properties
provided by the system.
In Section \ref{sec:delimitations}, we describe de-limitations of capabilitites
ofn the system.
In Section \ref{sec:services}, we describe the abstract service
functions consumed or produced by the system.
In Section \ref{sec:security}, we describe the security capabilitites
of the system.

\newpage

\subsection{Significant Prior Art}
\label{sec:prior_art}

Eclipse Arrowhead has its roots in service oriented arcitecture, SOA,
and its use for primarily far edge, edge and fog automation and
digitalsiation with interoperability to the cloud level.

A set of EU projects have built the foundation for whats now Eclipse
Arrowhead, current version 4.6.1 with the specifications for v5.0 in
the works. These projects are:
\begin{itemize}
\item Socrades
\item IMC-AESOP
\item Arrowhead
\item Productive4.0
\item Arrowhead Tools
\item AIMS5.0
\item Arrowhead fPVN
\end{itemize}

 The current comprehensive high level and architecture description of
 Eclipse Arrowhead architecture is the book ``IoT Automation -
 Arrowhead framework'' the  





\subsection{How This System Is Meant to Be Used}
\label{sec:use}

\color{red}
Describe intended usage of the system. Usage scenarios
shall be described. Preferable a SysML/UML blaock diagram of the
System should be provided. See the SysML profile and library
(github.com/eclipse-arrowhead/profile-library-sysml) for
support on how such block diagram should look like. Suitable tools are
Eclipse Papyrus and MagicDraw.
\color{black}  

\subsection{System functionalities and properties}
\label{sec:properties}

\color{red}
Narrative describe system functionalities and properties (no
implmentation details) like e.g.:

\subsubsection {Functional properties of the system}

\subsubsection {Configuration of system properties}


\subsubsection {Data stored by the system}
Brief overview of data stored to achive the functionality of the system. 

\subsubsection {Non functional properties}
\begin{itemize}
  \item security, 
  \item safety, 
  \item energy consumption,
  \item latency
  \item Power saving properties, 
\end{itemize}


\subsubsection {Stateful or stateless}
\begin{itemize} 
\item states preserved, functional and non-functional
\end{itemize}  
\color{black}  


\subsection{Important Delimitations}
\label{sec:delimitations}

\color{red}
Provide delimitations of the provided system. Describe what the system
solve and what i does not solve.
\color{black}  



\newpage

\section{Services}
\label{sec:services}

\color{red}
This section describes consumed and produced service.
In particular, each subsection names a prodiuced or consumed service
indicating the different capabilities and associated interfaces of the
service. Reference to the appropriate SD document shall be made.

\subsection{Produced service}
with references to SD and IDD documents

\subsection{Consumed services}
with references to SD and IDD documents

\color{black}




\newpage

\section{Security}
\label{sec:security}


\color{red}
Overview of security leel chosen for the system

The follwoing bullets should be covered 
\begin{itemize}
\item  If the system can be started in un-secure and/or
Arrowhead secure mode.
\item Handling of Arrowhead compliant and
non-compliant X.509 certificates.
\end{itemize}

\subsection {Security model}
The following points should be described:
\begin{itemize}
\item protocol supported 
\item data protection supported 
\item system authentication capability supported
\item produced service authorisation checking, 
\item etc.
\end{itemize}

For Arrowhead certificate profile
see github.com/eclipse-arrowhead/documentation
\color{black}




\bibliographystyle{IEEEtran}
\bibliography{bibliography}

\newpage

\section{Revision History}
\subsection{Amendments}

\color{red}
Revision history and Quality assurance as per examples below
\color{black}

\noindent\begin{tabularx}{\textwidth}{| p{1cm} | p{3cm} | p{2cm} | X | p{4cm} |} \hline
\rowcolor{gray!33} No. & Date & Version & Subject of Amendments & Author \\ \hline

1 & 2020-12-05 & \arrowversion & & Tanyi Szvetlin \\ \hline
2 & 2021-07-14 & \arrowversion & Minor updates & Jerker Delsing \\ \hline
3 & 2022-01-12 & \arrowversion & Minor updates & Jerker Delsing \\ \hline
\end{tabularx}

\subsection{Quality Assurance}

\noindent\begin{tabularx}{\textwidth}{| p{1cm} | p{3cm} | p{2cm} | X |} \hline
\rowcolor{gray!33} No. & Date & Version & Approved by \\ \hline

1 & 2022-01-10 & \arrowversion  &  \\ \hline

\end{tabularx}

\end{document}
\message{ !name(GSoSD_foundational principle developments.tex) !offset(-286) }
